% !TEX root = ../LAFROX-Subdocumento9-Anexos.tex
\chapter*{Anexos del Subdocumento 9: Plan de calidad}
\addcontentsline{toc}{chapter}{Anexos del Subdocumento 9: Plan de calidad}
\markboth{Anexos del Subdocumento 9: Plan de calidad}{}
\label{cap:sd9-anexos}

Estos anexos contienen el detalle que el cuerpo del Subdocumento 9 (SD9) resume. El plan de pruebas está en el Formulario T-13 y el protocolo de aceptación, en el Formulario T-17. En estos anexos, BA, BTT y Caso designan las Bases Administrativas, las Bases Técnicas Transversales y las Bases Técnicas del caso.

\begin{tablasimple}{>{\hsize=0.41\hsize\linewidth=\hsize}Y >{\hsize=2.18\hsize\linewidth=\hsize}Y >{\hsize=0.41\hsize\linewidth=\hsize}Y}
  {\cab{Anexo} & \cab{Contenido} & \cab{Sección del SD9 que lo cita}}
9.A & Métricas de calidad del producto por característica de ISO/IEC 25010 y de calidad de datos según ISO/IEC 25012 & 9.1.3 \\
9.B & Puertas de calidad, reglas de bloqueo y herramientas por lenguaje & 9.1.5 y 9.2.1 \\
9.C & Cálculo del costo de la calidad en horas-hombre & 9.1.6 y 9.3.1 \\
9.D & Trazabilidad de los resultados del caso hasta su prueba y su hito & 9.2.5 \\
\end{tablasimple}
\fuentesd{Fuente: (Subdocumento 9, secciones 9.1 a 9.3).}

\section*{Anexo 9.A — Métricas de calidad del producto}
\addcontentsline{toc}{section}{Anexo 9.A — Métricas de calidad del producto}

Las métricas siguen el modelo de calidad del producto de ISO/IEC 25010 (International Organization for Standardization [ISO] \& International Electrotechnical Commission [IEC], 2011) y el modelo de calidad de datos de ISO/IEC 25012 (ISO \& IEC, 2008). Cada métrica indica su umbral, su origen, el método con que se verifica y el paquete del Formulario T-14 donde se ejecuta. Los nombres de las características son los del Art. 24 de las Bases Administrativas (BA, p. 17).

\subsection*{9.A.1 Funcionalidad}
\addcontentsline{toc}{subsection}{9.A.1 Funcionalidad}

La funcionalidad se verifica contra los requisitos ofertados y los casos que acepta la Contraparte Técnica.

\begin{tablasimple}{>{\hsize=0.626\hsize\linewidth=\hsize}Y >{\hsize=0.626\hsize\linewidth=\hsize}Y >{\hsize=2.221\hsize\linewidth=\hsize}Y >{\hsize=0.626\hsize\linewidth=\hsize}Y >{\hsize=0.626\hsize\linewidth=\hsize}Y >{\hsize=1.652\hsize\linewidth=\hsize}Y >{\hsize=0.623\hsize\linewidth=\hsize}Y}
  {\cab{ID} & \cab{Subcaracterística} & \cab{Métrica} & \cab{Umbral} & \cab{Origen} & \cab{Verificación} & \cab{Paquete}}
9A-01 & Completitud & Requisitos ofertados del T-12 con prueba aprobada & 100 \% & Caso, 17.1 & Matriz de trazabilidad generada desde la cadena de integración & 3.8.1; 3.9.1 \\
9A-02 & Corrección & Defectos críticos o altos abiertos al certificar & 0 & Art. 17.3 (BA, p. 13) & Registro de defectos & 3.8.7; 3.9.6 \\
9A-03 & Pertinencia & Casos de aceptación firmados por la Contraparte Técnica & 100 \% de los casos de la etapa & BTT 20.1 (BTT, pp. 34--35) & Aceptación en Preproducción & 3.8.2; 3.9.2 \\
\end{tablasimple}
\fuentesd{Fuente: (Bases Administrativas, arts. 17.3 y 24, pp. 13, 17); (Bases Técnicas Transversales, sección 20.1, pp. 34--35); (Bases Técnicas del caso, sección 17.1, pp. 31--32); (Formulario T-13, sección 6); (Formulario T-14, cuentas 3 y 4).}

\subsection*{9.A.2 Desempeño}
\addcontentsline{toc}{subsection}{9.A.2 Desempeño}

El desempeño se mide en percentil 95 con la carga de peak y la concurrencia de la solución; no se sustituye por tiempos promedio.

\begin{tablasimple}{>{\hsize=0.575\hsize\linewidth=\hsize}Y >{\hsize=0.843\hsize\linewidth=\hsize}Y >{\hsize=1.717\hsize\linewidth=\hsize}Y >{\hsize=1.583\hsize\linewidth=\hsize}Y >{\hsize=0.851\hsize\linewidth=\hsize}Y >{\hsize=0.843\hsize\linewidth=\hsize}Y >{\hsize=0.588\hsize\linewidth=\hsize}Y}
  {\cab{ID} & \cab{Subcaracterística} & \cab{Métrica} & \cab{Umbral} & \cab{Origen} & \cab{Verificación} & \cab{Paquete}}
9A-04 & Comportamiento temporal & p95 de confirmación de línea de preparación & ≤ 1 s & Caso, RT-09.01 & Carga a 21,99 TPS & 3.8.4; 3.9.3 \\
9A-05 & Comportamiento temporal & p95 de registro de entrega & ≤ 2 s & Caso, RT-09.01 & Carga a 21,99 TPS & 3.8.4; 3.9.3 \\
9A-06 & Comportamiento temporal & p95 de línea de preventa & ≤ 1,5 s & Caso, RT-09.01 & Carga a 21,99 TPS & 3.8.4; 3.9.3 \\
9A-07 & Comportamiento temporal & p95 de consulta de stock y crédito & ≤ 2 s & Caso, RT-09.01 & Carga a 21,99 TPS & 3.8.4; 3.9.3 \\
9A-08 & Comportamiento temporal & p95 de transacción crítica de terreno, portal, navegación, API, búsqueda e informe & ≤ 3 s; ≤ 2 s; ≤ 1 s; consulta ≤ 500 ms; escritura ≤ 800 ms; búsqueda compuesta ≤ 3 s y global ≤ 2 s; informe estándar ≤ 30 s & BTT, cap. 9; SD4, sección 4.2.6.11 (BTT, pp. 21--22) & Carga a 21,99 TPS & 3.8.4; 3.9.3 \\
9A-09 & Capacidad & Carga sostenida sin violar 9A-04 a 9A-08 & 21,99 TPS; 1.163 personas usuarias; 237 dispositivos & Art. 24; RT-09.06; SD4, sección 4.2.6 (BA, p. 17; BTT, p. 21) & Carga en PRE-C & 3.8.4; 3.9.3 \\
9A-10 & Capacidad & Ventana de despacho a 1,5 × su peak & 10,41 TPS durante 90 minutos, 0 minutos de indisponibilidad & Caso, RT-10.05; SD4, sección 4.2.6 & Carga en PRE-C & 3.8.4; 3.9.3 \\
9A-11 & Capacidad & Crecimiento sin rediseño & 43,98 TPS (3 × el peak) sin cambio de arquitectura & RT-09.03 (BTT, p. 21) & Escalón de estrés & 3.8.4; 3.9.3 \\
9A-12 & Capacidad & Degradación al superar la capacidad & 0 transacciones perdidas; encolamiento y mensaje explícito & RT-09.08 (BTT, p. 22) & Escalón de quiebre (4 ×) & 3.8.4; 3.9.3 \\
\end{tablasimple}
\fuentesd{Fuente: (Bases Administrativas, art. 24, p. 17); (Bases Técnicas Transversales, cap. 9, pp. 21--22); (Bases Técnicas del caso, cap. 15, RT-09.01 y RT-10.05, pp. 26--27); (Subdocumento 4, secciones 4.2.6.3 y 4.2.6.4); (Formulario T-13, sección 7.1); (Formulario T-14, cuenta 3).}

\subsection*{9.A.3 Compatibilidad}
\addcontentsline{toc}{subsection}{9.A.3 Compatibilidad}

La compatibilidad se verifica sobre los quince flujos y sus contratos, incluida la convivencia con los sistemas del CLIENTE.

\begin{tablasimple}{>{\hsize=0.654\hsize\linewidth=\hsize}Y >{\hsize=0.835\hsize\linewidth=\hsize}Y >{\hsize=2.038\hsize\linewidth=\hsize}Y >{\hsize=0.654\hsize\linewidth=\hsize}Y >{\hsize=1.242\hsize\linewidth=\hsize}Y >{\hsize=0.924\hsize\linewidth=\hsize}Y >{\hsize=0.653\hsize\linewidth=\hsize}Y}
  {\cab{ID} & \cab{Subcaracterística} & \cab{Métrica} & \cab{Umbral} & \cab{Origen} & \cab{Verificación} & \cab{Paquete}}
9A-13 & Interoperabilidad & Flujos de integración ejecutados sin error por etapa & 100 \% de E1; 15 de 15 al cierre de E2 & BTT 20.1; SD4, sección 4.1 (BTT, pp. 34--35) & Integración en QA & 3.8.1; 3.9.1 \\
9A-14 & Interoperabilidad & Contratos rotos en la cadena de integración & 0 & Art. 4.3 (OpenAPI 3.1, AsyncAPI 2.6) (BA, pp. 5--6) & Pact y Spectral & 1.5.2 \\
9A-15 & Coexistencia & Regresiones de la Etapa 1 al integrar la Etapa 2 & 0 & Art. 17.2, regla 3 (BA, p. 13) & Regresión de la Etapa 1 & 3.9.1 \\
\end{tablasimple}
\fuentesd{Fuente: (Bases Técnicas Transversales, sección 20.1 y RT-10.08, pp. 22, 34--35); (Subdocumento 4, sección 4.1); (Formulario T-13, sección 3); (Formulario T-14, cuentas 1 y 3).}

\subsection*{9.A.4 Usabilidad}
\addcontentsline{toc}{subsection}{9.A.4 Usabilidad}

La usabilidad combina accesibilidad y resultados medidos por perfil en los dispositivos y condiciones reales de trabajo.

\begin{tablasimple}{>{\hsize=0.514\hsize\linewidth=\hsize}Y >{\hsize=0.88\hsize\linewidth=\hsize}Y >{\hsize=2.319\hsize\linewidth=\hsize}Y >{\hsize=1.136\hsize\linewidth=\hsize}Y >{\hsize=0.514\hsize\linewidth=\hsize}Y >{\hsize=1.123\hsize\linewidth=\hsize}Y >{\hsize=0.514\hsize\linewidth=\hsize}Y}
  {\cab{ID} & \cab{Subcaracterística} & \cab{Métrica} & \cab{Umbral} & \cab{Origen} & \cab{Verificación} & \cab{Paquete}}
9A-16 & Accesibilidad & Incumplimientos de WCAG 2.2 A y AA & 0, con informe de conformidad & RT-13.01; Art. 4.3 (BA, pp. 5--6; BTT, p. 26) & axe-core y revisión manual & 3.8.2; 3.9.2 \\
9A-17 & Operabilidad & Condiciones de terreno superadas: guantes a \ensuremath{-}22 °C, lluvia, sobre 34 °C, sol directo, una mano & 5 de 5 & Caso, RT-13.08 & Prueba de las cinco condiciones físicas, separada de 9A-25; regresión del alcance cambiado E2 & 3.8.3; 3.9.1; 3.9.2 \\
9A-18 & Eficiencia de uso y aprendizaje & Tiempo, pasos, tasa de error y tiempo de aprendizaje por perfil & Preparador ≤ 1 s, 3 pasos, error ≤ 5 \%, aprendizaje ≤ 2 h; recepción ≤ 2 s, 3 pasos, ≤ 5 \%, ≤ 2 h; preventista ≤ 1,5 s, 3 pasos, ≤ 5 \%, ≤ 4 h; conductor ≤ 2 s, 3 o 4 pasos, ≤ 5 \%, ≤ 30 min & RT-13.04; SD3, Tabla 3.A.4a (BTT, p. 26) & Pruebas moderadas por perfil & 2.6.3; 3.8.2 \\
9A-19 & Adaptabilidad de la interfaz & Funciones completas en escritorio, tableta y teléfono & 100 \% de las vistas de portales y consolas & RT-13.02 (BTT, p. 26) & Regresión visual con Playwright & 3.8.1; 3.9.1 \\
\end{tablasimple}
\fuentesd{Fuente: (Bases Administrativas, art. 4.3, pp. 5--6); (Bases Técnicas Transversales, RT-13.01 a RT-13.04, p. 26); (Bases Técnicas del caso, cap. 15, RT-13.08, p. 27); (Subdocumento 3, Tabla 3.A.4a); (Formulario T-13, secciones 6 y 7.3); (Formulario T-14, cuentas 2 y 3).}

\subsection*{9.A.5 Fiabilidad}
\addcontentsline{toc}{subsection}{9.A.5 Fiabilidad}

La fiabilidad se verifica mediante disponibilidad, recuperación e inyección de fallas, con umbrales por nivel de servicio.

\begin{tablasimple}{>{\hsize=0.498\hsize\linewidth=\hsize}Y >{\hsize=0.706\hsize\linewidth=\hsize}Y >{\hsize=1.942\hsize\linewidth=\hsize}Y >{\hsize=1.304\hsize\linewidth=\hsize}Y >{\hsize=0.651\hsize\linewidth=\hsize}Y >{\hsize=1.192\hsize\linewidth=\hsize}Y >{\hsize=0.707\hsize\linewidth=\hsize}Y}
  {\cab{ID} & \cab{Subcaracterística} & \cab{Métrica} & \cab{Umbral} & \cab{Origen} & \cab{Verificación} & \cab{Paquete}}
9A-20 & Disponibilidad & Disponibilidad mensual por nivel de servicio y de la infraestructura & Crítico ≥ 99,9 \%; alto ≥ 99,5 \%; medio ≥ 99,0 \%; bajo ≥ 98,0 \%; infraestructura (energía, climatización, red, cómputo, base de datos y portal) ≥ 99,95 \% & Art. 78.2; BTT 7.2; SD4, sección 4.2.4.3 (BA, p. 40; BTT, pp. 17--18) & Observabilidad en marcha blanca y Operación & 4.2.1; 4.3.1; Operación \\
9A-21 & Disponibilidad & Indisponibilidad entre 05:30 y 07:00, lunes a sábado & 0 minutos & Caso, RT-10.05 & Indicador diario de la marcha blanca & 4.2.1; 4.3.1 \\
9A-22 & Tolerancia a fallos & Escenarios de inyección superados sin intervención & 8 de 8 (T-13, sección 7.2) & RT-10.07; SD4, Anexo 4-V (BTT, p. 22) & AWS FIS, sitio emulado y gateways de frío & 3.8.4; 3.9.3; 8.2.2 \\
9A-23 & Recuperabilidad & RTO y RPO en conmutación real & RTO ≤ 4 h; RPO ≤ 15 min & RT-07.04; RT-07.07 (BTT, p. 17) & Conmutación regional y recuperación del WMS de Talca en la nube (AL-DR-01) & 3.8.5; 3.9.4; 8.1.3 \\
9A-24 & Recuperabilidad & Restauraciones mensuales exitosas, con tiempo medido & 100 \% & RT-07.12 (BTT, p. 18) & Restauración de muestra & 4.2.1; 4.2.2; 4.3.1; 8.1.6 \\
9A-25 & Operación desconectada & Registros perdidos o duplicados tras 14 h sin señal y 24 h sin enlace & 0; sincronización ≤ 10 min (dispositivo) y ≤ 2 h (centro de distribución) & Caso, RT-03.10 y RT-03.13 & Perfil operacional y resiliencia; regresión del alcance cambiado E2 & 3.8.3; 3.8.4; 3.9.1; 3.9.2 \\
\end{tablasimple}
\fuentesd{Fuente: (Bases Administrativas, art. 78.2, p. 40); (Bases Técnicas Transversales, sección 7.2, RT-07.07 y RT-10.07, pp. 17--18, 22); (Bases Técnicas del caso, cap. 15, RT-10.05, p. 27); (Formulario T-13, secciones 7.2 y 9); (Formulario T-14, cuentas 3 y 8).}

\subsection*{9.A.6 Seguridad}
\addcontentsline{toc}{subsection}{9.A.6 Seguridad}

La seguridad se evalúa por hallazgos abiertos y por la aplicación efectiva de los controles de identidad y datos.

\begin{tablasimple}{>{\hsize=0.546\hsize\linewidth=\hsize}Y >{\hsize=0.637\hsize\linewidth=\hsize}Y >{\hsize=2.047\hsize\linewidth=\hsize}Y >{\hsize=0.77\hsize\linewidth=\hsize}Y >{\hsize=0.959\hsize\linewidth=\hsize}Y >{\hsize=1.462\hsize\linewidth=\hsize}Y >{\hsize=0.579\hsize\linewidth=\hsize}Y}
  {\cab{ID} & \cab{Subcaracterística} & \cab{Métrica} & \cab{Umbral} & \cab{Origen} & \cab{Verificación} & \cab{Paquete}}
9A-26 & Resistencia & Hallazgos críticos o altos abiertos (SAST, SCA, DAST, imágenes, intrusión) & 0 & RT-04.05; RT-11.20 (BTT, pp. 10, 24) & Cadena de integración; tercero independiente & 1.5.2; 3.8.6; 3.9.5; 8.2.2 \\
9A-27 & Resistencia & Plazo de corrección de hallazgos & Crítico 7 días corridos; alto 15; medio 30 & BA, Art. 21.1; SD4, sección 4.2.4.1.4 (BA, p. 15) & Registro de vulnerabilidades & 8.2.2 \\
9A-28 & Confidencialidad & Datos productivos sin anonimizar fuera de Producción & 0; 0 hallazgos de Amazon Macie & Ley N.º 21.719; Caso, RT-11.10 & Verificación previa a la liberación de cada juego & Todas las campañas \\
9A-29 & Confidencialidad & Secretos en código, imágenes o configuración & 0 & RT-04.09 (BTT, p. 11) & Escaneo de secretos & 1.5.2 \\
\end{tablasimple}
\fuentesd{Fuente: (Bases Técnicas Transversales, RT-04.05 y RT-11.20, pp. 10, 24); (Bases Técnicas del caso, cap. 15, pp. 26--27); (Subdocumento 4, sección 4.1.18); (Formulario T-13, secciones 3 y 6); (Formulario T-14, cuentas 1 y 3).}

\subsection*{9.A.7 Mantenibilidad}
\addcontentsline{toc}{subsection}{9.A.7 Mantenibilidad}

La mantenibilidad se controla antes de fusionar código. Cada métrica debe producir evidencia verificable y bloquear el cambio que exceda su umbral.

\begin{tablasimple}{>{\hsize=0.694\hsize\linewidth=\hsize}Y >{\hsize=0.822\hsize\linewidth=\hsize}Y >{\hsize=1.76\hsize\linewidth=\hsize}Y >{\hsize=0.694\hsize\linewidth=\hsize}Y >{\hsize=1.187\hsize\linewidth=\hsize}Y >{\hsize=1.147\hsize\linewidth=\hsize}Y >{\hsize=0.696\hsize\linewidth=\hsize}Y}
  {\cab{ID} & \cab{Subcaracterística} & \cab{Métrica} & \cab{Umbral} & \cab{Origen} & \cab{Verificación} & \cab{Paquete}}
9A-30 & Testeabilidad & Cobertura de la lógica de negocio & ≥ 70 \% & Art. 24; RT-04.11 (BA, p. 17; BTT, p. 11) & PHPUnit con PCOV; Kover; Jest & 1.5.2 \\
9A-31 & Testeabilidad & Cobertura unitaria global y de líneas modificadas, medidas por separado & ≥ 80 \% en cada denominador & Política de LafroX (SD1, 1.3.1); T-10, DevSecOps & PHPUnit con PCOV; Kover; Jest, con dos informes separados & 1.5.2 \\
9A-32 & Analizabilidad & Complejidad ciclomática por método & ≤ 10 & McCabe (1976) & PHPMD; detekt; ESLint & 1.5.2 \\
9A-33 & Modularidad & Dependencias prohibidas entre módulos; acoplamiento eferente por clase & 0; ≤ 20 clases & Art. 19 (acoplamiento débil); regla de LafroX (BA, p. 14) & deptrac; PhpMetrics; detekt; ESLint & 1.5.2 \\
9A-34 & Modificabilidad & Duplicación de código & ≤ 3 \% de las líneas & Formulario T-10 & PHPCPD; jscpd & 1.5.2 \\
9A-35 & Modificabilidad & Deuda técnica clasificada como bloqueante & 0 & Art. 24; Formulario T-10 (BA, p. 17) & PhpMetrics y registro de deuda & 1.5.2; 8.2.4 \\
\end{tablasimple}
\fuentesd{Fuente: (Bases Administrativas, arts. 19 y 24, pp. 14, 17); (Bases Técnicas Transversales, RT-04.11, p. 11); (Subdocumento 1, sección 1.3.1); (Formulario T-10, apartado DevSecOps); McCabe (1976, sección III); (Formulario T-14, cuentas 1 y 8).}

\subsection*{9.A.8 Portabilidad}
\addcontentsline{toc}{subsection}{9.A.8 Portabilidad}

La portabilidad se comprueba promoviendo el mismo artefacto y reconstruyendo el ambiente desde sus definiciones versionadas.

\begin{tablasimple}{>{\hsize=0.598\hsize\linewidth=\hsize}Y >{\hsize=0.723\hsize\linewidth=\hsize}Y >{\hsize=1.891\hsize\linewidth=\hsize}Y >{\hsize=0.783\hsize\linewidth=\hsize}Y >{\hsize=0.598\hsize\linewidth=\hsize}Y >{\hsize=1.81\hsize\linewidth=\hsize}Y >{\hsize=0.597\hsize\linewidth=\hsize}Y}
  {\cab{ID} & \cab{Subcaracterística} & \cab{Métrica} & \cab{Umbral} & \cab{Origen} & \cab{Verificación} & \cab{Paquete}}
9A-36 & Instalabilidad & Despliegues que promueven el mismo artefacto sin recompilar & 100 \% & RT-04.08 (BTT, p. 11) & Digest de la imagen en cada ambiente & 3.8.8; 3.9.7 \\
9A-37 & Adaptabilidad & Ambientes reconstruibles desde código & 5 de 5 & BTT 4.1 (BTT, p. 10) & Inventario y código de los cinco ambientes; reconstrucción aislada de cada configuración con datos sintéticos y evidencia de equivalencia. Producción/Recuperación se validan sin destruir instancias activas & 3.1 \\
9A-38 & Reemplazabilidad & Despliegue y reversión sin interrupción & 0 interrupciones; 0 intervenciones manuales & RT-04.06; RT-04.07 (BTT, p. 11) & Demostración en Preproducción & 3.8.8; 3.9.7 \\
\end{tablasimple}
\fuentesd{Fuente: (Bases Técnicas Transversales, RT-04.06 a RT-04.08 y sección 4.1, pp. 10--11); (Subdocumento 4, sección 4.2.4.1); (Formulario T-13, secciones 4 y 6); (Formulario T-14, cuentas 3 y 4).}

\subsection*{9.A.9 Calidad de datos (ISO/IEC 25012)}
\addcontentsline{toc}{subsection}{9.A.9 Calidad de datos (ISO/IEC 25012)}

La calidad de datos se mide sobre registros del negocio. Se distinguen completitud, exactitud, consistencia y actualidad para no ocultar errores con un único indicador.

\begin{tablasimple}{>{\hsize=0.588\hsize\linewidth=\hsize}Y >{\hsize=0.588\hsize\linewidth=\hsize}Y >{\hsize=2.524\hsize\linewidth=\hsize}Y >{\hsize=0.588\hsize\linewidth=\hsize}Y >{\hsize=0.958\hsize\linewidth=\hsize}Y >{\hsize=1.167\hsize\linewidth=\hsize}Y >{\hsize=0.587\hsize\linewidth=\hsize}Y}
  {\cab{ID} & \cab{Dimensión} & \cab{Métrica} & \cab{Umbral} & \cab{Origen} & \cab{Verificación} & \cab{Paquete}}
9A-39 & Completitud & Recepciones con lote registrado, en productos que lo requieren & 100 \% & Caso, cap. 18, resultado 2 & Indicador de la marcha blanca & 4.2.1 \\
9A-40 & Exactitud & Diferencias no explicadas en la conciliación de la migración & 0 & BTT 20.1; SD5, 5.3.5 (BTT, pp. 34--35) & Dos ensayos y corte & 3.7.2 a 3.7.5 \\
9A-41 & Consistencia & Definiciones del indicador de entregas completas y a tiempo & 1 sola & Caso, cap. 18, resultado 4 & Revisión del modelo analítico & 4.2.1 \\
9A-42 & Actualidad & Latencia de los indicadores de la operación del día & ≤ 5 min & Caso, RT-05.29 & Medición en marcha blanca & 4.2.1 \\
\end{tablasimple}
\fuentesd{Fuente: (Bases Técnicas del caso, cap. 18, pp. 34--35); (Bases Técnicas Transversales, RT-04.04 y sección 20.1, pp. 10, 34--35); (Subdocumento 5, sección 5.3.5); (Formulario T-13, secciones 4 y 7.4); (Formulario T-14, cuentas 3 y 4).}

\section*{Anexo 9.B — Puertas de calidad y reglas de bloqueo}
\addcontentsline{toc}{section}{Anexo 9.B — Puertas de calidad y reglas de bloqueo}

Las reglas siguientes concretan qué impide una promoción en cada puerta y qué evidencia se exige para corregirlo; ninguna aprobación sustituye una prueba fallida.

\begin{tablasimple}{>{\hsize=0.325\hsize\linewidth=\hsize}Y >{\hsize=0.544\hsize\linewidth=\hsize}Y >{\hsize=1.775\hsize\linewidth=\hsize}Y >{\hsize=1.418\hsize\linewidth=\hsize}Y >{\hsize=0.938\hsize\linewidth=\hsize}Y}
  {\cab{Puerta} & \cab{Momento} & \cab{Reglas que bloquean} & \cab{Herramienta (PHP · Kotlin · TypeScript)} & \cab{Quién puede levantar el bloqueo}}
G0 & Compromiso local & Formato; pruebas del módulo en falla & Laravel Pint · ktlint · ESLint; pruebas unitarias & Nadie; se corrige \\
G1 & Solicitud de fusión & Falta la aprobación de un par; secreto detectado; error de lint & Revisión en GitLab; escaneo de secretos & Nadie \\
G2 & Integración continua & 9A-30 a 9A-35 incumplidos; hallazgo crítico o alto de SAST o SCA; contrato roto (9A-14); prueba en falla & PHPStan con Larastan, PHPMD, PHPCPD, PhpMetrics, deptrac · detekt, Kover, jscpd · ESLint, Jest; SAST de GitLab; \texttt{composer audit}; Pact; Spectral & Nadie \\
G3 & Promoción a QA & Imagen con vulnerabilidad crítica o alta; regresión de desempeño: p95, errores o colas de la carga de regresión peores que en la última versión aceptada, bajo la misma carga y configuración (SD4, sección 4.2.4.1.1) & Amazon Inspector; k6 & Nadie; se corrige \\
G4 & Promoción a Preproducción & Regresión con fallas; alerta alta de DAST o MobSF; incumplimiento de WCAG 2.2 AA; defecto crítico o alto abierto & Playwright; Espresso; OWASP ZAP; MobSF; axe-core & Nadie \\
G5 & Entrada de cada etapa a Producción (H6 y H11) & Una prueba del T-13 que habilita la entrada no está aprobada; defecto crítico o alto abierto; falta el acta de certificación de la etapa (H5 o H10) & Expediente del T-17 & Nadie. El acta es una condición adicional y no levanta una prueba fallida \\
G6 & Producción desde H6/H11 y Operación & Canario con p95 sobre su umbral o con más errores que la versión estable; presupuesto de error crítico del mes agotado & Cadena de entrega del SD4; observabilidad & Nadie. La reversión del canario es automática; se suspenden los despliegues no correctivos hasta el mes siguiente \\
\end{tablasimple}
\fuentesd{Fuente: (Bases Técnicas Transversales, RT-04.05 y RT-04.11, pp. 10--11); (Subdocumento 4, sección 4.2.4.1); (Formulario T-10, apartado DevSecOps); (Formulario T-13, sección 6); (Formulario T-17, sección 4.3).}

G5 es la puerta de entrada de cada etapa a Producción en el H6 y el H11 (T-17, 4.3): bloquea la habilitación sin certificación H5/H10 aceptada, pruebas aplicables aprobadas, reversión ensayada, conciliación y capacitación. El H7 y el H12 no son una puerta técnica: son el paso a producción contractual del Formulario E-25 y aceptan la marcha blanca con las seis condiciones del Art. 17.3 (BA, p. 13).

Los umbrales de las puertas son \textbf{tolerancias}: el rango que un requisito acepta, como el p95 de 1 s de la línea de preparación. Son distintos de los \textbf{límites de control} de la sección 9.2.6 del SD9, que se calculan con la variación histórica de una métrica y sirven para detectar un proceso inestable antes de que llegue a la tolerancia (PMI, 2017, pp. 274 y 304).

\section*{Anexo 9.C — Costo de la calidad en horas-hombre}
\addcontentsline{toc}{section}{Anexo 9.C — Costo de la calidad en horas-hombre}

El costo de la calidad se expresa en horas-hombre y no en dinero, porque la Oferta Técnica no puede contener cifras que permitan inferir el monto de la oferta (BA, Art. 50.2, p. 29). Las categorías son las del PMBOK: prevención, evaluación y fallas (PMI, 2017, pp. 274 y 282–283).

\subsection*{9.C.1 Costos de conformidad (programados en la línea base)}
\addcontentsline{toc}{subsection}{9.C.1 Costos de conformidad (programados en la línea base)}

El subtotal reúne los paquetes de prevención y evaluación identificables en la línea base. Cada paquete se cuenta una sola vez. Banco 2.4.3, mapeo 2.4.4 y piloto 4.1.4 se clasifican como evaluación porque producen evidencia de conformidad física y aceptación de la cadena de frío: 320 + 240 + 320 = 880 HH. La clase de estimación E del mapeo no determina su categoría de costo de calidad. Estas horas ya pertenecen a las 220.023 HH del T-15; no son un incremento de esa línea base ni se vuelven a contar en las pruebas de integración.

\begin{tablasimple}{>{\hsize=0.729\hsize\linewidth=\hsize}Y >{\hsize=0.664\hsize\linewidth=\hsize}Y >{\hsize=1.943\hsize\linewidth=\hsize}Y >{\hsize=0.664\hsize\linewidth=\hsize}Y}
  {\cab{Categoría} & \cab{Paquete T-14} & \cab{Contenido} & \cab{HH (T-15)}}
Prevención & 1.2.4 & Matriz de trazabilidad & 80 \\
Prevención & 1.5.1 & Plan de calidad y de pruebas & 80 \\
Prevención & 1.5.2 & Puertas de calidad automáticas & 80 \\
Prevención & 2.6.1 a 2.6.3 & Investigación con usuarios, prototipos e indicadores de usabilidad & 720 \\
Prevención & 3.10.2.1 a 3.10.2.4 & INN-02: reproducción de incidentes de terreno & 960 \\
Prevención & 8.2.4 & Registro y reducción de la deuda técnica en la Operación & 2.304 \\
\textbf{Subtotal prevención} &  &  & \textbf{4.224} \\
Evaluación & 2.4.3; 2.4.4; 4.1.4 & Banco (320), mapeo validado por Calidad (240) y piloto de frío (320) & 880 \\
Evaluación & 3.8.1 a 3.8.8 & Pruebas y certificación de la Etapa 1 & 2.240 \\
Evaluación & 3.9.1 a 3.9.7 & Pruebas y certificación de la Etapa 2 & 1.600 \\
Evaluación & 4.1.2 & Ensayo del procedimiento de reversión & 80 \\
Evaluación & 8.1.3 & Pruebas de recuperación ante desastres en la Operación & 560 \\
Evaluación & 8.1.6 & Restauraciones mensuales en la Operación & 576 \\
Evaluación & 8.2.2 & Vulnerabilidades, intrusión anual e inyección de fallas & 1.152 \\
Evaluación & 8.3.1 & Regresión de INN-02 en la Operación & 576 \\
Evaluación & 4.2.1; 4.2.2; 4.3.1 & Ocho restauraciones previas a Operación, componentes SRE adicionales & 128 \\
\textbf{Subtotal evaluación} &  &  & \textbf{7.792} \\
\textbf{Total de conformidad} &  &  & \textbf{12.016} \\
\end{tablasimple}
\fuentesd{Fuente: (Formulario T-14, paquetes indicados en cada fila); (Formulario T-15, estimación por paquete y rol).}

Las 12.016 HH equivalen al 5,5 \% de las 220.023 HH de la línea base (12.016 / 220.023 = 0,0546).

\textbf{Partidas excluidas.} No se cuentan la migración de datos (cuenta 3.7), cuyos ensayos forman parte del trabajo de migración y no se pueden separar de él, ni la certificación de los usuarios del CLIENTE (7.3.1 y 7.3.2), que verifica personas y no el producto.

\subsection*{9.C.2 Fallas y autorización de recursos}
\addcontentsline{toc}{subsection}{9.C.2 Fallas y autorización de recursos}

El costo de no conformidad se estima por evento y por especialidad, sin usar importes esperados de un registro de riesgos. Se distingue falla interna, detectada por LafroX antes de entregar, de falla externa, detectada en el uso del CLIENTE. Las cantidades siguientes son dos escenarios de referencia de LafroX para planificar retrabajo: no son mediciones, frecuencia prevista ni reserva autorizada.

En ambos escenarios se supone una corrección acotada que exige un ciclo manual completo de 270 casos a 15 minutos: 270 × 15 / 60 = 67,50 HH, presentado como 67,5 HH. La parte automática de la regresión dura aproximadamente 9,0 horas transcurridas en un ejecutor; ese tiempo no se suma como HH de una persona. Las HH de preparación, supervisión y análisis de ese ciclo están incluidas en SRE y CAL. Si una corrección modifica la arquitectura, los datos o la superficie de seguridad, se reestima también la campaña especializada afectada.

\begin{tablasimple}{Y Y Y Y}
  {\cab{Rol} & \cab{Falla interna, HH/evento} & \cab{Falla externa, HH/evento} & \cab{Desglose de referencia}}
Desarrollo (DES) & 16 & 24 & Interna: 4 análisis + 8 corrección + 4 pruebas unitarias; externa: 8 + 12 + 4 \\
Calidad (CAL) & 75,5 & 75,5 & 67,5 de ejecución manual + 4 de análisis + 4 de evidencia \\
Operación (SRE) & 4 & 8 & Interna: 2 de restitución + 2 de supervisión; externa: 4 de diagnóstico + 4 de promoción y verificación \\
Implantación y coordinación (IMP) & 0 & 4 & Comunicación y coordinación con las personas afectadas por la falla externa \\
\textbf{Total de referencia} & \textbf{95,5} & \textbf{111,5} & HH por evento; no constituyen dotación ni autorización de gasto \\
\end{tablasimple}
\fuentesd{Fuente: (Formulario T-13, secciones 8.2 y 8.3); (Formulario T-15, estimación por rol); los supuestos de esfuerzo por evento explicitados en esta sección 9.C.2.}

Para nI fallas internas y nE externas de este alcance, el esfuerzo de referencia es 95,5 × nI + 111,5 × nE. Un evento de cada tipo representa 207 HH. El costo de calidad del escenario sería 12.016 + 207 = 12.223 HH únicamente si esas horas fueran adicionales y no estuvieran ya incluidas en los paquetes de la línea base. No se adopta ese escenario como pronóstico del contrato: las cantidades nI y nE se revisan con los defectos observados desde las primeras campañas.

Cada evento registra su requisito, severidad, etapa, causa, versión, HH previstas y reales por rol, pruebas repetidas y paquete de imputación. Se diferencia la repetición causada por el defecto de las pruebas planificadas que ya pertenecen a evaluación. Antes de agregar esfuerzo a la línea base se descuentan las HH ya incluidas y asignadas al mismo trabajo; no se vuelve a sumar la capacidad protegida del T-15 ni la atención correctiva incluida en 8.2.1. El Líder de Calidad informa mensualmente costos internos y externos consumidos y el Jefe de Proyecto autoriza los recursos por control de cambios tras verificar capacidad, sin consumir cobertura de servicio ni reservas protegidas de otro frente.

\subsection*{9.C.3 Distribución del esfuerzo programado}
\addcontentsline{toc}{subsection}{9.C.3 Distribución del esfuerzo programado}

Las proporciones siguientes comparan prevención y evaluación dentro del subtotal de conformidad; no representan la distribución del costo integral con fallas.

\begin{tablasimple}{>{\hsize=1.052\hsize\linewidth=\hsize}Y >{\hsize=0.974\hsize\linewidth=\hsize}Y >{\hsize=0.974\hsize\linewidth=\hsize}Y}
  {\cab{Categoría} & \cab{HH} & \cab{Proporción}}
Prevención & 4.224 & 35,2 \% \\
Evaluación & 7.792 & 64,8 \% \\
\textbf{Total programado} & \textbf{12.016} & \textbf{100,0 \%} \\
\end{tablasimple}
\fuentesd{Fuente: (Formulario T-14, paquetes de prevención y evaluación); (Formulario T-15, estimación por paquete y rol); el desglose de la sección 9.C.1.}

Proporciones: 4.224 / 12.016 = 0,3515; 7.792 / 12.016 = 0,6485. El T-15 desglosa campañas por actividad y rol: recuperación inicial 320 HH; siete conmutaciones operacionales de 80 HH (560); intrusión E1/E2 de 320/160 HH. Son estimaciones de planificación, no resultados ejecutados ni valores derivados sólo de la proporción de alcance.

\section*{Anexo 9.D — Trazabilidad de los resultados del caso}
\addcontentsline{toc}{section}{Anexo 9.D — Trazabilidad de los resultados del caso}

La matriz prevista individual está publicada en T-13, sección 8.2, y se indexa en SD9, sección 9.2.5. El entregable 1.2.4 (H1) la completa con las fichas ejecutables y con el registro de evidencias al ejecutar; no parte de un conteo agregado sin identificadores. Su actualización se genera desde el repositorio: cada caso de prueba lleva el identificador de su requisito, cada solicitud de fusión cita el requisito o la incidencia que atiende y cada despliegue registra el digest de su imagen (RT-04.04). En la oferta, la correspondencia entre requisito, componente y sección está en el Formulario T-12.

Este anexo muestra el tramo final de la cadena del Caso, 17.1, para los dieciséis resultados de aceptación: de cada resultado a la prueba que lo ensaya antes de la marcha blanca, al paquete donde se mide y al hito que lo acepta.

\begin{tablasimple}{>{\hsize=1.461\hsize\linewidth=\hsize}Y >{\hsize=1.55\hsize\linewidth=\hsize}Y >{\hsize=0.495\hsize\linewidth=\hsize}Y >{\hsize=0.494\hsize\linewidth=\hsize}Y}
  {\cab{Resultado} & \cab{Prueba previa (T-13)} & \cab{Paquete de medición} & \cab{Hito de aceptación}}
R18-01 Retiro sanitario en < 2 h & Simulacro sobre datos sintéticos (sección 7.4) & 4.2.1 & H7 \\
R18-02 100 \% de recepciones con lote & Casos de recepción; 9A-39 & 4.2.1 & H7 \\
R18-03 Registro continuo de temperatura & Banco 2.4.3, mapeo 2.4.4 y piloto 4.1.4; V1–V21 del Anexo 4-W, excursiones y retención por SSCC (T-13, sección 7.4) & 4.2.1 & H7 \\
R18-04 OTIF medido de una forma & Revisión del modelo analítico; 9A-41 & 4.2.1, 4.3.1 y Operación & Por tramo (meses 15, 19 y 32) \\
R18-05 Stock y crédito en la preventa & 9A-07 & 4.2.1 & H7 \\
R18-06 Ningún pedido perdido sin señal & Jornada de 14 h sin señal (sección 7.3); 9A-25 & 4.2.1 & H7 \\
R18-07 Ruta en < 20 min & Generación con volumen de peak (sección 7.4) & 4.2.1 & H7 \\
R18-08 Prueba de entrega digital el mismo día & Casos de entrega en aceptación & 4.2.1 & H7 \\
R18-09 Cero guías extraviadas o ilegibles & Casos de emisión y consulta de guías & 4.2.1 & H7 \\
R18-10 Rendición cuadrada el mismo día & Casos de cierre por camión & 4.2.1 & H7 \\
R18-11 Costo de servir por cliente & Casos de la Etapa 2 en aceptación & 4.3.1 & H12 \\
R18-12 Control de envases & Casos de saldo por cliente y transportista & 4.2.1 y Operación & H7 (provisional); tras 12 meses de operación, mes 32 (final) \\
R18-13 Pedidos EDI de cadenas & Contratos de integración EDI; 9A-13 & 4.3.1 & H12 \\
R18-14 Ocupación mínima por camión & Casos de planificación de carga & 4.2.1 & H7 \\
R18-15 Causa de cada faltante & Casos de faltante sin conexión & 4.2.1 & H7 \\
R18-16 Operación sin el planificador & Generación automática de rutas (R18-07) & 4.2.1 & H7 \\
\end{tablasimple}
\fuentesd{Fuente: (Bases Técnicas del caso, cap. 18, pp. 34--35); (Formulario T-13, sección 7.4); (Formulario T-14, cuentas 3 y 4); (Formulario T-17, secciones 4 y 5).}

\section*{Referencias}
\addcontentsline{toc}{section}{Referencias}

Las fuentes siguientes sustentan los criterios citados en este documento; sus referencias se leen junto con el artículo, requisito o sección indicado en el lugar de uso.

\begin{itemize}[before=\raggedright,itemsep=0pt,parsep=0pt,topsep=2pt]
  \item Distribuidora Puelche S.A. (2026a). \emph{Bases Administrativas, Licitación TFEP-01/2026}.
  \item Distribuidora Puelche S.A. (2026b). \emph{Bases Técnicas Transversales, Licitación TFEP-01/2026}.
  \item Distribuidora Puelche S.A. (2026c). \emph{Bases Técnicas del Caso Logística}.
  \item International Organization for Standardization \& International Electrotechnical Commission. (2008). \emph{ISO/IEC 25012:2008. Software engineering — Software product Quality Requirements and Evaluation (SQuaRE) — Data quality model}.
  \item International Organization for Standardization \& International Electrotechnical Commission. (2011). \emph{ISO/IEC 25010:2011. Systems and software engineering — Systems and software Quality Requirements and Evaluation (SQuaRE) — System and software quality models}.
  \item Ley N.º 21.719. (2024). \emph{Regula la protección y el tratamiento de los datos personales y crea la Agencia de Protección de Datos Personales}. Diario Oficial de la República de Chile.
  \item McCabe, T. J. (1976). A complexity measure. \emph{IEEE Transactions on Software Engineering, SE-2}(4), 308–320.
  \item Project Management Institute. (2017). \emph{Guía de los fundamentos para la dirección de proyectos (Guía del PMBOK)} (6.ª ed.). Project Management Institute.
\end{itemize}
